% 第1章 1.6.1：闭子空间上的正交投影（示意图）。平面表示闭子空间 M，Pf 为最近点，
% 残差 f-Pf 与 M 垂直；对任意 m∈M，||f-m||^2=||f-Pf||^2+||Pf-m||^2。
\begin{tikzpicture}[font=\figurefont\footnotesize,line cap=round]
\fill[FigureBlueFill] (0,0) -- (5.6,0) -- (7.0,1.9) -- (1.4,1.9) -- cycle;
\draw[FigureBlue,line width=.75pt] (0,0) -- (5.6,0) -- (7.0,1.9) -- (1.4,1.9) -- cycle;
\node[text=FigureBlue,anchor=south west] at (0.25,0.05) {$M$};
\coordinate (O) at (1.6,0.55);
\coordinate (P) at (4.0,1.0);
\coordinate (F) at (4.0,3.6);
\coordinate (Q) at (5.6,1.45);
\draw[FigureStroke,line width=.6pt,-{Stealth[length=2mm]}] (O) -- (F);
\draw[FigureBlue,line width=1.1pt,-{Stealth[length=2mm]}] (O) -- (P);
\draw[FigureYellow,line width=1.1pt] (P) -- (F);
\draw[FigureGreen,line width=.9pt,dashed] (F) -- (Q);
\draw[FigureGreen,line width=.9pt] (P) -- (Q);
\draw[FigureStroke,line width=.5pt] ($(P)+(0,0.28)$) -- ($(P)+(0.27,0.33)$) -- ($(P)+(0.27,0.05)$);
\fill[FigureStroke] (O) circle (1.2pt) node[anchor=north east,text=FigureInk] {$0$};
\fill[FigureStroke] (P) circle (1.4pt) node[anchor=north,text=FigureInk,yshift=-1pt] {$Pf$};
\fill[FigureStroke] (F) circle (1.4pt) node[anchor=south,text=FigureInk] {$f$};
\fill[FigureStroke] (Q) circle (1.4pt) node[anchor=west,text=FigureInk] {$m$};
\node[text=FigureYellow,anchor=east] at ($(P)!0.5!(F)$) {$f-Pf$};
\node[text=FigureGreen,anchor=south west] at ($(F)!0.55!(Q)$) {$f-m$};
\end{tikzpicture}
